\documentclass{beamer}

\usepackage{amsmath,amsfonts,amssymb}
\usepackage{physics}
\usepackage{bm}
\usepackage{quantikz}

\title{Circuit Depth and Trotter Time Steps in Variational Quantum Algorithms}
\subtitle{Application to the Unitary Coupled Cluster Ansatz}
\author{MHJ}
\date{}

\begin{document}

%------------------------------------------------
\frame{\titlepage}

%------------------------------------------------
\begin{frame}{Motivation}

Variational quantum eigensolvers (VQE) rely on preparing
parameterized quantum states

\[
|\Psi(\theta)\rangle = U(\theta)|\Phi_0\rangle
\]

and minimizing the energy

\[
E(\theta) = \langle \Psi(\theta)|H|\Psi(\theta)\rangle.
\]

Many physically motivated ansätze are generated by
exponentials of operators,

\[
U(\theta) = e^{-iA}.
\]

Examples:

\begin{itemize}
\item Hamiltonian evolution
\item Unitary coupled cluster
\item time-evolution inspired ansätze
\end{itemize}

\end{frame}

%------------------------------------------------
\begin{frame}{Operator Exponentials}

Consider a unitary generated by a sum of operators

\[
U = e^{-iA}
\]

where

\[
A = \sum_k A_k .
\]

In general

\[
[A_i , A_j] \neq 0.
\]

Therefore the exponential cannot be implemented directly
as a product of exponentials.

\end{frame}

%------------------------------------------------
\begin{frame}{Suzuki–Trotter Decomposition}

To implement the exponential we use the Trotter formula

\[
e^{-iA} =
e^{-i\sum_k A_k}
\approx
\left(
\prod_k e^{-i A_k \Delta t}
\right)^n
\]

where

\[
n\Delta t = 1.
\]

Thus

\[
\Delta t = \frac{1}{n}.
\]

\end{frame}

%------------------------------------------------
\begin{frame}{Interpretation}

The unitary operator is approximated by

\[
U(\Delta t) =
\prod_k e^{-i A_k \Delta t}.
\]

To obtain the full transformation we apply

\[
n = \frac{1}{\Delta t}
\]

Trotter steps.

\end{frame}

%------------------------------------------------
\begin{frame}{Circuit Depth}

Each operator

\[
e^{-iA_k\Delta t}
\]

corresponds to a small quantum circuit.

If there are $M$ operators

\[
A = \sum_{k=1}^M A_k
\]

then one Trotter step contains $M$ circuit blocks.

\end{frame}

%------------------------------------------------
\begin{frame}{Total Circuit Length}

The full circuit contains

\[
N_{\text{gates}}
\sim
n \times M .
\]

Using

\[
n = \frac{1}{\Delta t}
\]

we obtain

\[
N_{\text{gates}} \sim \frac{M}{\Delta t}.
\]

\end{frame}

%------------------------------------------------
\begin{frame}{Key Result}

\[
\boxed{
\text{Circuit depth} \propto \frac{1}{\Delta t}
}
\]

Small Trotter step

\[
\Delta t \rightarrow 0
\]

improves accuracy but increases circuit length.

\end{frame}

%------------------------------------------------
\begin{frame}{Trotter Error}

The leading error term is

\[
\mathcal{O}(\Delta t^2)
\]

coming from commutators

\[
[A_i , A_j].
\]

Thus accuracy improves when

\[
\Delta t \rightarrow 0.
\]

\end{frame}

%------------------------------------------------
\begin{frame}{Tradeoff}

There is a fundamental tradeoff

\begin{itemize}
\item smaller $\Delta t$ to  higher accuracy
\item smaller $\Delta t$ to  deeper circuits
\end{itemize}

This tradeoff is central in variational quantum algorithms.

\end{frame}

%------------------------------------------------
\section{Application to Unitary Coupled Cluster}

%------------------------------------------------
\begin{frame}{Unitary Coupled Cluster Ansatz}

The UCC wavefunction is

\[
|\Psi_{UCC}\rangle =
e^{T-T^\dagger}|\Phi_0\rangle .
\]

Define

\[
\tau = T-T^\dagger .
\]

The unitary operator becomes

\[
U = e^{\tau}.
\]

\end{frame}

%------------------------------------------------
\begin{frame}{Cluster Operator}

The doubles operator is

\[
T_2 =
\frac{1}{4}
\sum_{abij}
t_{ij}^{ab}
a_a^\dagger a_b^\dagger a_j a_i .
\]

Thus

\[
\tau =
\sum_k \tau_k
\]

where each $\tau_k$ represents a fermionic excitation.

\end{frame}

%------------------------------------------------
\begin{frame}{UCC as Sum of Generators}

The UCC generator is

\[
\tau =
\sum_k \theta_k G_k
\]

with

\[
G_k =
a_a^\dagger a_b^\dagger a_j a_i
-
a_i^\dagger a_j^\dagger a_b a_a .
\]

These operators do not commute.

\end{frame}

%------------------------------------------------
\begin{frame}{Trotterized UCC Ansatz}

The exponential becomes

\[
e^{\sum_k \theta_k G_k}
\approx
\left(
\prod_k
e^{\theta_k G_k \Delta t}
\right)^n .
\]

\end{frame}

%------------------------------------------------
\begin{frame}{Number of Excitations}

For $N$ orbitals

\[
N_{\text{exc}} \sim N^4 .
\]

Thus one Trotter step contains

\[
N^4
\]

excitation circuits.

\end{frame}

%------------------------------------------------
\begin{frame}{Total Circuit Depth}

Using

\[
n = \frac{1}{\Delta t}
\]

the circuit depth becomes

\[
D \sim \frac{N^4}{\Delta t}.
\]

This scaling explains why UCC circuits quickly become large.

\end{frame}


\begin{frame}{Example Circuit Block}

After Jordan--Wigner mapping the excitation becomes a Pauli rotation.

\begin{center}
\begin{quantikz}
\lstick{$q_0$} & \ctrl{1} & \qw      & \qw      & \qw \\
\lstick{$q_1$} & \targ{}  & \ctrl{1} & \qw      & \qw \\
\lstick{$q_2$} & \qw      & \targ{}  & \ctrl{1} & \qw \\
\lstick{$q_3$} & \qw      & \qw      & \targ{}  & \gate{$R_z(\theta)$}
\end{quantikz}
\end{center}

\end{frame}

%------------------------------------------------
\begin{frame}{Depth Scaling}

Each excitation requires

\begin{itemize}
\item CNOT ladder
\item single qubit rotation
\item inverse ladder
\end{itemize}

Thus the total gate count grows rapidly with

\[
\frac{1}{\Delta t}.
\]

\end{frame}

%------------------------------------------------
\begin{frame}{Implications}

Because

\[
D \propto \frac{1}{\Delta t}
\]

UCC circuits become deep when high Trotter accuracy is required.

Strategies to reduce circuit depth include

\begin{itemize}
\item operator grouping
\item adaptive ansätze
\item parametric matrix models
\end{itemize}

\end{frame}

%------------------------------------------------
\begin{frame}{Summary}

\begin{itemize}
\item Many variational ansätze involve exponentials of operator sums
\item Noncommuting generators require Trotter decomposition
\item Number of Trotter steps

\[
n = \frac{1}{\Delta t}
\]

\item Circuit depth therefore scales as

\[
D \propto \frac{1}{\Delta t}.
\]

\item For UCC with $N$ orbitals

\[
D \sim \frac{N^4}{\Delta t}.
\]

\end{itemize}

\end{frame}

\end{document}
